\section{SCAN/r2SCAN meta-GGA energies and gradients}
\label{sec:mgga}

\subsection{Semilocal ingredients}

At each grid point, a spin-unpolarized MGGA evaluates
\begin{equation}
  e_{\mathrm{xc}}
  =f(\rho,\sigma,\tau),
  \qquad \sigma=\nabla\rho\cdot\nabla\rho.
  \label{eq:mgga-ingredients}
\end{equation}
The weights entering the variational derivative are
\begin{equation}
  v_\rho=\frac{\partial f}{\partial\rho},
  \qquad
  v_\sigma=\frac{\partial f}{\partial\sigma},
  \qquad
  v_\tau=\frac{\partial f}{\partial\tau}.
  \label{eq:mgga-weights}
\end{equation}

SCAN is a constraint-based functional designed to connect the uniform-electron
gas, single-orbital, and slowly varying density limits.  r2SCAN regularizes the
numerically sensitive switching functions of SCAN and improves grid
convergence, but it has the same derivative structure in terms of $\rho$,
$\sigma$, and $\tau$.

\subsection{AO representation of \texorpdfstring{$\tau$}{tau}}

Writing Eq.~\eqref{eq:tau-definition} as an AO product gives
\begin{equation}
  \tau(\mathbf r)
  =\frac{1}{2N_k}\sum_{\vk\mu\nu}D_{\mu\nu}^{\vk}
   \sum_{y=x,y,z}
   \partial_y\phi_{\nu\vk}^{*}(\mathbf r)
   \partial_y\phi_{\mu\vk}(\mathbf r).
  \label{eq:mgga-tau-ao}
\end{equation}
The tau contribution to the target MGGA Fock matrix is
\begin{equation}
  V_{\mu\nu}^{\tau,\vk}
  =\frac12\int_\Omega v_\tau(\mathbf r)
   \nabla\phi_{\nu\vk}^{*}(\mathbf r)\cdot
   \nabla\phi_{\mu\vk}(\mathbf r)\dd\mathbf r.
  \label{eq:mgga-vtau-matrix}
\end{equation}
The factor $1/2$ comes from Eq.~\eqref{eq:tau-definition}.  One must first
determine whether the numerical integration library already incorporates the
$\tau$ convention into its effective weight and then apply the factor exactly
once.

\subsection{Explicit nuclear derivative}

At fixed density matrix,
\begin{equation}
  E_{\mathrm{xc}}^{Ax}
  =\int_\Omega\left(
    v_\rho\rho^{Ax}
    +v_\sigma\sigma^{Ax}
    +v_\tau\tau^{Ax}\right)\dd\mathbf r.
  \label{eq:mgga-xc-derivative}
\end{equation}
On a fixed-cell uniform grid, quadrature points and volume weights do not
change with a nuclear coordinate, so only AO derivatives remain.  The GGA
part contains
\begin{equation}
  \sigma^{Ax}=2\nabla\rho\cdot\nabla\rho^{Ax},
  \label{eq:mgga-sigma-derivative}
\end{equation}
while the tau derivative is
\begin{align}
  \tau^{Ax}(\mathbf r)
  =\frac{1}{2N_k}\sum_{\vk\mu\nu}D_{\mu\nu}^{\vk}
   \sum_y\big[&
   \partial_y\phi_{\nu\vk}^{Ax *}
   \partial_y\phi_{\mu\vk}
  \notag\\
   &+\partial_y\phi_{\nu\vk}^{*}
   \partial_y\phi_{\mu\vk}^{Ax}\big].
  \label{eq:mgga-tau-derivative}
\end{align}
Because an AO-center derivative and a spatial gradient appear together,
second spatial derivatives of the AOs are required.  This is the principal
reason why an MGGA gradient kernel is more complicated than a GGA kernel.

\subsection{Coupling to the nonvariational response}

HF still defines the source orbitals in an MGGA target calculation.  The
Hessian $A$ is therefore the same as in Eq.~\eqref{eq:hf-hessian-action}.  The
addition of Eq.~\eqref{eq:mgga-vtau-matrix} to the target Fock matrix changes
the target VO gradient to
\begin{equation}
  g_{ai}^{\vk,\mathrm{MGGA}}
  =\left[C_v^{\vk\dagger}
   \left(h+J+V_\rho+V_\sigma+V_\tau\right)^{\vk}
   C_o^{\vk}\right]_{ai}.
  \label{eq:mgga-target-vo}
\end{equation}
After solving for the Z-vector, it is contracted with the source-HF nuclear
perturbation through Eq.~\eqref{eq:response-gradient}.  Connecting the MGGA
kernel only to the explicit XC term while retaining a PBE target-Fock VO
right-hand side would give an incorrect nonvariational response.

\subsection{Grid aliasing}

High-frequency components of Gaussian AOs alias on a coarse grid.  Since
$\tau$ is a product of AO gradients, it is more sensitive to high reciprocal
components than the density itself.  MGGA calculations therefore require a
separate convergence study of \texttt{cell.mesh}, or of the effective
integration mesh.  The resolved integer triplet and cutoff should be recorded
together rather than changing one direction at a time.

The translational sum rule is a sensitive measure of grid aliasing.  Even when
the analytic--FD component error is small, a large
\begin{equation}
  \left\|\sum_A\mathbf g_A^{\mathrm{xc}}\right\|
  \label{eq:mgga-xc-translation}
\end{equation}
means that the explicit XC-grid derivative has not converged.

\subsection{CPU/GPU parity and backend composition}

CPU and GPU routes must use the same libxc functional label, density
convention, mesh, and AO derivative order.  Comparisons include not only total
energy and final gradient, but also $E_{\mathrm{xc}}$, the target-Fock VO
block, the $v_\tau$ row derivative, the Z-vector right-hand side, and its
residual.  Completion cases for SCAN/r2SCAN include Gamma-point H$_2$,
distorted H$_2$O, and sampled k-points.

The target MGGA and the source-HF density-fitting provider have separate
responsibilities.  The same $\rho$, $\nabla\rho$, and $\tau$ target path must
therefore be tested with RSDF as well as FFTDF and GDF.  An RSDF+MGGA case is a
composition test in which the XC derivative of
Eq.~\eqref{eq:mgga-xc-derivative} and the source/metric auxiliary response of
Eq.~\eqref{eq:rsdf-total-gradient} operate simultaneously.
